% 编译：xelatex -halt-on-error assignment7-sol.tex（运行两遍）。
% 中文使用 TeX Live 自带的 Fandol 字体。
\documentclass[11pt,a4paper,fontset=fandol]{ctexart}
\usepackage[margin=2.05cm]{geometry}
\usepackage{amsmath,amssymb,mathtools}
\usepackage{fancyhdr,needspace}
\usepackage[hidelinks]{hyperref}
\hypersetup{pdftitle={线性代数第七次作业},pdfsubject={MATH1205：矩阵的秩、初等变换与矩阵方程}}
\setlength{\parindent}{0pt}
\setlength{\parskip}{0.2em}
\setlength{\headheight}{15pt}
\renewcommand{\arraystretch}{1.08}
\linespread{1.15}
\pagestyle{fancy}
\fancyhf{}
\fancyhead[L]{MATH1205\quad 线性代数}
\fancyhead[R]{第七次作业}
\fancyfoot[C]{\thepage}
\renewcommand{\headrulewidth}{0.3pt}
\ctexset{section={format=\large\bfseries,beforeskip=1.2ex,afterskip=0.5ex},
subsection={format=\normalsize\bfseries,beforeskip=0.9ex,afterskip=0.4ex}}
\newcommand{\rk}{\operatorname{rank}}
\newcommand{\ans}[1]{\boxed{#1}}
\newcommand{\problem}[1]{\Needspace{8\baselineskip}\section*{#1}}
\newcommand{\parttitle}[1]{\Needspace{7\baselineskip}\subsection*{#1}}

\begin{document}
\setlength{\abovedisplayskip}{6pt plus 2pt minus 2pt}
\setlength{\belowdisplayskip}{6pt plus 2pt minus 2pt}
\setlength{\abovedisplayshortskip}{3pt plus 1pt minus 1pt}
\setlength{\belowdisplayshortskip}{4pt plus 1pt minus 1pt}
\begin{center}
{\LARGE\bfseries 线性代数第七次作业}\par
\vspace{0.3em}
{\small MATH1205}
\end{center}
\vspace{-0.4em}
题号沿用原题。$E_k$ 表示 $k$ 阶单位矩阵，$r_i$ 表示矩阵的第 $i$ 行。

\problem{第 45 题：参数与矩阵的秩}
题中
\[
A=\begin{pmatrix}
1&a&a&a\\a&1&a&a\\a&a&1&a\\a&a&a&1
\end{pmatrix}.
\]
在 $|A|$ 中把其余各列加到第一列，提出第一列的公因子 $1+3a$，
再把第一行分别从其余三行中减去，得
\[
|A|=(1+3a)
\begin{vmatrix}
1&a&a&a\\1&1&a&a\\1&a&1&a\\1&a&a&1
\end{vmatrix}
=(1+3a)
\begin{vmatrix}
1&a&a&a\\0&1-a&0&0\\0&0&1-a&0\\0&0&0&1-a
\end{vmatrix}
=(1+3a)(1-a)^3.
\]
\textbf{(1)} $A$ 可逆当且仅当 $|A|\ne0$，故
\[
\ans{a\ne1\quad\text{且}\quad a\ne-\frac13}.
\]
\textbf{(2)} 秩为 $3$ 必有 $|A|=0$，只需考察 $a=1,-\frac13$。
当 $a=1$ 时，$A$ 是全 $1$ 矩阵，秩为 $1$；当 $a=-\frac13$ 时，
原矩阵前 $3$ 行、前 $3$ 列构成的子式为
\[
\begin{vmatrix}
1&-\frac13&-\frac13\\
-\frac13&1&-\frac13\\
-\frac13&-\frac13&1
\end{vmatrix}=\frac{16}{27}\ne0.
\]
因此此时 $\rk A=3$，所以 $\ans{a=-\frac13}$。

\textbf{(3)} 由以上讨论，其余参数值的秩均为 $4$ 或 $3$，故
$\ans{\rk A=1\iff a=1}$。

\problem{第 47 题：用定义求秩}
\parttitle{(3)}
设
\[
B=\begin{pmatrix}
1&2&3&4&5\\0&0&-1&-2&-3\\0&0&0&0&4\\0&0&1&2&-1
\end{pmatrix}.
\]
取第 $1,2,3$ 行与第 $1,3,5$ 列，得到三阶子式
\[
\begin{vmatrix}1&3&5\\0&-1&-3\\0&0&4\end{vmatrix}=-4\ne0.
\]
又因 $r_2+r_3+r_4=0$，任取四列所得的四阶子式中，仍有上述行关系，
故所有四阶子式均为零；更高阶子式不存在。由秩的定义，$\ans{\rk B=3}$。

\newpage
\section*{第 47 题（续）}
\parttitle{(4)}
设 $C=(a_i b_j)_{n\times n}$。任意二阶子式均有
\[
\begin{vmatrix}a_i b_j&a_i b_\ell\\a_k b_j&a_k b_\ell\end{vmatrix}
=a_i a_k b_j b_\ell-a_i a_k b_\ell b_j=0.
\]
因而 $\rk C\le1$（当 $n=1$ 时也显然成立）。若 $a_i$ 全为零或 $b_j$ 全为零，
则 $C=O$，秩为 $0$；否则存在 $a_i\ne0$ 及 $b_j\ne0$，一阶子式 $a_i b_j\ne0$。
所以
\[
\ans{\rk C=
\begin{cases}
0,&(a_1,\ldots,a_n)=0\ \text{或}\ (b_1,\ldots,b_n)=0,\\
1,&(a_1,\ldots,a_n)\ne0\ \text{且}\ (b_1,\ldots,b_n)\ne0.
\end{cases}}
\]

\problem{第 48 题：初等变换求秩及最高阶非零子式}
\parttitle{(1)}
对原矩阵作行变换：
\[
A_1=\begin{pmatrix}3&1&0&2\\1&-1&2&-1\\1&3&-4&4\end{pmatrix}
\xrightarrow{r_1\leftrightarrow r_2}
\begin{pmatrix}1&-1&2&-1\\3&1&0&2\\1&3&-4&4\end{pmatrix}
\xrightarrow{\substack{r_2\leftarrow r_2-3r_1\\r_3\leftarrow r_3-r_1}}
\begin{pmatrix}1&-1&2&-1\\0&4&-6&5\\0&4&-6&5\end{pmatrix}.
\]
再作 $r_3\leftarrow r_3-r_2$，得阶梯形矩阵
\[
\begin{pmatrix}1&-1&2&-1\\0&4&-6&5\\0&0&0&0\end{pmatrix},
\qquad \ans{\rk A_1=2}.
\]
原矩阵第 $1,2$ 行与第 $1,2$ 列构成的子式为
\[
\ans{\begin{vmatrix}3&1\\1&-1\end{vmatrix}=-4\ne0}.
\]

\parttitle{(3)}
先对原矩阵作 $r_1\leftrightarrow r_4$，再作
$r_2\leftarrow r_2-2r_1$、$r_3\leftarrow r_3-3r_1$、$r_4\leftarrow r_4-2r_1$，得
\[
A_3=\begin{pmatrix}
2&1&8&3&7\\2&-3&0&7&-5\\3&-2&5&8&0\\1&0&3&2&0
\end{pmatrix}
\longrightarrow
\begin{pmatrix}
1&0&3&2&0\\0&-3&-6&3&-5\\0&-2&-4&2&0\\0&1&2&-1&7
\end{pmatrix}.
\]
交换 $r_2,r_4$，再作 $r_3\leftarrow r_3+2r_2$、$r_4\leftarrow r_4+3r_2$，随后
$r_4\leftarrow r_4-\frac87r_3$，得到
\[
\begin{pmatrix}
1&0&3&2&0\\0&1&2&-1&7\\0&0&0&0&14\\0&0&0&0&16
\end{pmatrix}
\longrightarrow
\begin{pmatrix}
1&0&3&2&0\\0&1&2&-1&7\\0&0&0&0&14\\0&0&0&0&0
\end{pmatrix}.
\]
故 $\ans{\rk A_3=3}$。原矩阵第 $1,2,4$ 行与第 $1,2,5$ 列构成的子式为
\[
\ans{\begin{vmatrix}2&1&7\\2&-3&-5\\1&0&0\end{vmatrix}
=\begin{vmatrix}1&7\\-3&-5\end{vmatrix}=16\ne0}.
\]

\newpage
\section*{第 48 题（续）}
\parttitle{(5)}
对原矩阵作 $r_2\leftarrow r_2-2r_1$、$r_3\leftarrow r_3-3r_1$、
$r_4\leftarrow r_4-r_1$，得到
\[
A_5=\begin{pmatrix}1&3&-1&2\\2&-1&2&3\\3&2&1&1\\1&-4&3&5\end{pmatrix}
\longrightarrow
\begin{pmatrix}1&3&-1&2\\0&-7&4&-1\\0&-7&4&-5\\0&-7&4&3\end{pmatrix}.
\]
再作 $r_3\leftarrow r_3-r_2$、$r_4\leftarrow r_4-r_2$，随后 $r_4\leftarrow r_4+r_3$，得
\[
\begin{pmatrix}1&3&-1&2\\0&-7&4&-1\\0&0&0&-4\\0&0&0&4\end{pmatrix}
\longrightarrow
\begin{pmatrix}1&3&-1&2\\0&-7&4&-1\\0&0&0&-4\\0&0&0&0\end{pmatrix}.
\]
故 $\ans{\rk A_5=3}$。原矩阵第 $1,2,3$ 行与第 $1,2,4$ 列构成的子式为
\[
\ans{\begin{vmatrix}1&3&2\\2&-1&3\\3&2&1\end{vmatrix}
=\begin{vmatrix}1&3&2\\0&-7&-1\\0&0&-4\end{vmatrix}=28\ne0}.
\]

\problem{第 49 题：用初等变换求逆矩阵}
\parttitle{(2)}
记原矩阵为 $M$，对增广矩阵的左右两部分同时作初等行变换：
\[
(M\mid E_4)=\left(\begin{array}{rrrr|rrrr}
3&-2&0&-1&1&0&0&0\\
0&2&2&1&0&1&0&0\\
1&-2&-3&-2&0&0&1&0\\
0&1&2&1&0&0&0&1
\end{array}\right).
\]
先交换 $r_1,r_3$，并交换 $r_2,r_4$，再作 $r_3\leftarrow r_3-3r_1$。
接着作 $r_1\leftarrow r_1+2r_2$、$r_3\leftarrow r_3-4r_2$、$r_4\leftarrow r_4-2r_2$，得
\[
\left(\begin{array}{rrrr|rrrr}
1&0&1&0&0&0&1&2\\
0&1&2&1&0&0&0&1\\
0&0&1&1&1&0&-3&-4\\
0&0&-2&-1&0&1&0&-2
\end{array}\right)
\xrightarrow{r_4\leftarrow r_4+2r_3}
\left(\begin{array}{rrrr|rrrr}
1&0&1&0&0&0&1&2\\
0&1&2&1&0&0&0&1\\
0&0&1&1&1&0&-3&-4\\
0&0&0&1&2&1&-6&-10
\end{array}\right).
\]
最后依次作
\[
r_3\leftarrow r_3-r_4,\quad
r_1\leftarrow r_1-r_3,\quad
r_2\leftarrow r_2-2r_3,\quad
r_2\leftarrow r_2-r_4,
\]
左侧化为 $E_4$，右侧即为 $M^{-1}$。因此
\[
\ans{M^{-1}=\begin{pmatrix}
1&1&-2&-4\\0&1&0&-1\\-1&-1&3&6\\2&1&-6&-10
\end{pmatrix}}.
\]

\newpage
\section*{第 49 题（续）}
\parttitle{(4)}
记原矩阵为 $N$。对 $(N\mid E_4)$ 按 $i=4,3,2$ 的顺序作
$r_i\leftarrow r_i-r_{i-1}$，得
\[
\left(\begin{array}{rrrr|rrrr}
1&0&0&0&1&0&0&0\\
2&1&0&0&0&1&0&0\\
3&2&1&0&0&0&1&0\\
4&3&2&1&0&0&0&1
\end{array}\right)
\longrightarrow
\left(\begin{array}{rrrr|rrrr}
1&0&0&0&1&0&0&0\\
1&1&0&0&-1&1&0&0\\
1&1&1&0&0&-1&1&0\\
1&1&1&1&0&0&-1&1
\end{array}\right).
\]
再按 $i=4,3,2$ 的顺序作一次相同的变换，得到
\[
\left(\begin{array}{rrrr|rrrr}
1&0&0&0&1&0&0&0\\
0&1&0&0&-2&1&0&0\\
0&0&1&0&1&-2&1&0\\
0&0&0&1&0&1&-2&1
\end{array}\right).
\]
故
\[
\ans{N^{-1}=\begin{pmatrix}
1&0&0&0\\-2&1&0&0\\1&-2&1&0\\0&1&-2&1
\end{pmatrix}}.
\]

\problem{第 50 题：解矩阵方程}
\parttitle{(1)}
记
\[
P=\begin{pmatrix}2&5\\1&3\end{pmatrix},\qquad
Q=\begin{pmatrix}4&-6\\2&1\end{pmatrix}.
\]
因 $|P|=6-5=1\ne0$，故
\[
P^{-1}=\begin{pmatrix}3&-5\\-1&2\end{pmatrix}.
\]
由 $PX=Q$，在等式两端左乘 $P^{-1}$，得到唯一解
\[
\ans{X=P^{-1}Q
=\begin{pmatrix}3&-5\\-1&2\end{pmatrix}
\begin{pmatrix}4&-6\\2&1\end{pmatrix}
=\begin{pmatrix}2&-23\\0&8\end{pmatrix}}.
\]

\newpage
\section*{第 50 题（续）}
\parttitle{(3)}
题中
\[
A=\begin{pmatrix}3&0&1\\1&1&1\\1&1&4\end{pmatrix},\qquad
B=\begin{pmatrix}2&1&3\\0&1&2\\1&0&3\end{pmatrix}.
\]
将 $AX=B+2X$ 移项，得到
\[
(A-2E_3)X=B,\qquad
C:=A-2E_3=\begin{pmatrix}1&0&1\\1&-1&1\\1&1&2\end{pmatrix}.
\]
对 $(C\mid B)$ 作行变换：
\[
\left(\begin{array}{rrr|rrr}
1&0&1&2&1&3\\1&-1&1&0&1&2\\1&1&2&1&0&3
\end{array}\right)
\xrightarrow{\substack{r_2\leftarrow r_2-r_1\\r_3\leftarrow r_3-r_1}}
\left(\begin{array}{rrr|rrr}
1&0&1&2&1&3\\0&-1&0&-2&0&-1\\0&1&1&-1&-1&0
\end{array}\right).
\]
再依次作 $r_2\leftarrow-r_2$、$r_3\leftarrow r_3-r_2$，得
\[
\left(\begin{array}{rrr|rrr}
1&0&1&2&1&3\\0&1&0&2&0&1\\0&0&1&-3&-1&-1
\end{array}\right)
\xrightarrow{r_1\leftarrow r_1-r_3}
\left(\begin{array}{rrr|rrr}
1&0&0&5&2&4\\0&1&0&2&0&1\\0&0&1&-3&-1&-1
\end{array}\right).
\]
左侧已化为 $E_3$，说明 $C$ 可逆，原方程的唯一解为
\[
\ans{X=\begin{pmatrix}5&2&4\\2&0&1\\-3&-1&-1\end{pmatrix}}.
\]
代入检验：
\[
(A-2E_3)X
=\begin{pmatrix}1&0&1\\1&-1&1\\1&1&2\end{pmatrix}
\begin{pmatrix}5&2&4\\2&0&1\\-3&-1&-1\end{pmatrix}
=\begin{pmatrix}2&1&3\\0&1&2\\1&0&3\end{pmatrix}=B.
\]

\end{document}
